\documentclass[11pt,a4paper]{article}
\usepackage[margin=27mm]{geometry}\usepackage{amsmath,amssymb,amsthm}
\newtheorem{lemma}{Lemma}
\title{Two endpoint anchors on seven points}\author{Human Compression project}\date{}
\begin{document}\maketitle
\begin{lemma}
Let $a,b,c$ be involutions on a set $X$ of seven points, with $ac=ca$.
Suppose that distinct points $u,v\in X$ satisfy
\[
 a(u)\ne u,\quad b(u)=c(u)=u,
 \qquad c(v)\ne v,\quad a(v)=b(v)=v.
\]
Then $\langle a,b,c\rangle$ is intransitive on $X$.
\end{lemma}
\begin{proof}
Put $K=\langle a,c\rangle$. Since $a$ and $c$ commute and have order
two, every $K$-orbit has size $1$, $2$ or $4$. The orbit of $u$ is
exactly $\{u,a(u)\}$: indeed, $c$ fixes both points because
$c(a(u))=a(c(u))=a(u)$. Similarly, the orbit of $v$ is exactly
$\{v,c(v)\}$. These two orbits are distinct. Otherwise $v$ would
belong to the first orbit, on which $a$ interchanges its two points,
contrary to $a(v)=v$.

These orbits account for four points. The other three cannot form a
single $K$-orbit, since three is not an allowed orbit size. Hence $K$
has at least four orbits on $X$.

The involution $b$ fixes $u$ and $v$. It is therefore a product of at
most two disjoint transpositions on the other five points. Start with
the graph whose connected components are the $K$-orbits, and add an
edge for each transposition of $b$. Each added edge reduces the number
of components by at most one. There remain at least $4-2=2$ components.
They are precisely the orbits of $\langle a,b,c\rangle$, since paths
using the edges of the three involutions describe their generated orbits.
Thus the generated group is intransitive.
\end{proof}

\paragraph{Application and scope.}
A rank-three string representation on seven points with anchors at its
first and last marks satisfies the hypotheses. It cannot generate the
natural symmetric group. The conclusion requires neither the intersection
property nor two eligible edges for each mark. It resolves only this
seven-point boundary; it makes no assertion about higher degrees.

\paragraph{Verification and attribution.}
The written proof is self-contained. A separate enumeration normalizes
$u=1,v=7$ by relabelling. It examines all $231$ nonidentity involutions,
then all $200$ commuting endpoint pairs and $5000$ anchored triples.
None is transitive. An independent GAP enumeration using actual generated
group orbits gives the same counts. Positive controls in degrees four
and eight, and a seven-point control missing one anchor condition, show
that these tests do not merely reject every input. These computations
supplement the proof; they are not its completeness premise.
The requested boundary was specified by the investigator of the critical
endpoint argument. The orbit proof and the two enumerations were developed
here. No priority, whole-manuscript, conjecture, Lean or compression claim
is made.
\end{document}
